% Mathématiques 4e — cours V3
% Puissances et notation scientifique

\section{Objectifs}

\begin{itemize}
\item Comprendre une puissance comme écriture condensée d'un produit.
\item Utiliser des puissances de base quelconque à exposant entier positif.
\item Utiliser les puissances de $10$ à exposants positifs et négatifs.
\item Relier puissances de $10$ et préfixes d'unités de nano à giga.
\item Écrire un nombre en notation scientifique.
\item Passer de la notation scientifique à l'écriture décimale.
\item Comparer des nombres et déterminer des ordres de grandeur.
\item Effectuer des calculs en revenant à la définition des puissances.
\end{itemize}

\section{1. Puissance d'exposant positif}

Pour un nombre $a$ et un entier positif $n$ :
\[
a^n
\]
désigne un produit de $n$ facteurs égaux à $a$.

Par exemple :
\[
3^4=3\times3\times3\times3=81.
\]

\begin{definition}
Pour $n\ge1$ :
\[
\boxed{a^n=\text{produit de }n\text{ facteurs égaux à }a}.
\]
\end{definition}

\section{2. Exposant zéro}

Pour $a\neq0$ :
\[
\boxed{a^0=1}.
\]

En particulier :
\[
10^0=1.
\]

\section{3. Puissances de dix positives}

\[
10^1=10,\qquad 10^2=100,\qquad10^3=1000.
\]

Plus généralement :
\[
10^n
\]
est égal à $1$ suivi de $n$ zéros lorsque $n$ est positif.

\section{4. Puissances de dix négatives}

\begin{definition}
Pour $n>0$ :
\[
\boxed{10^{-n}=\dfrac1{10^n}}.
\]
\end{definition}

Ainsi :
\[
10^{-1}=0{,}1,
\]
\[
10^{-2}=0{,}01,
\]
\[
10^{-3}=0{,}001.
\]

\begin{attention}
Le signe négatif de l'exposant ne rend pas le nombre négatif.
\end{attention}

\section{5. Multiplier par une puissance de dix}

\[
4{,}27\times10^3=4270.
\]

Et :
\[
4{,}27\times10^{-3}=0{,}00427.
\]

Le sens du déplacement de la virgule peut être contrôlé par l'ordre de grandeur.

\section{6. Préfixes scientifiques}

\coursefigure{/images/mathematiques/4eme/puissances-prefixes.svg}{Préfixes de nano à giga associés à des puissances de dix.}{Les préfixes scientifiques permettent d'exprimer efficacement des grandeurs très petites ou très grandes.}

Quelques correspondances :
\[
\text{nano}=10^{-9},\quad
\text{micro}=10^{-6},\quad
\text{milli}=10^{-3},
\]
\[
\text{kilo}=10^3,\quad
\text{méga}=10^6,\quad
\text{giga}=10^9.
\]

\section{7. Conversion avec un préfixe}

\[
1\ \text{nm}=10^{-9}\ \text{m}.
\]

Donc :
\[
450\ \text{nm}=450\times10^{-9}\ \text{m}.
\]

Comme :
\[
450=4{,}5\times10^2,
\]
on obtient :
\[
450\ \text{nm}=4{,}5\times10^{-7}\ \text{m}.
\]

\section{8. Notation scientifique}

\coursefigure{/images/mathematiques/4eme/puissances-scientifique.svg}{Forme générale d'une notation scientifique.}{Le coefficient possède exactement un chiffre non nul avant la virgule.}

\begin{definition}
La notation scientifique d'un nombre non nul est :
\[
\boxed{a\times10^n}
\]
avec :
\[
1\le |a|<10
\]
et $n$ entier.
\end{definition}

\section{9. Grand nombre}

\[
45\,700\,000
=
4{,}57\times10^7.
\]

Le coefficient :
\[
4{,}57
\]
est bien compris entre $1$ et $10$.

\section{10. Petit nombre}

\[
0{,}000031
=
3{,}1\times10^{-5}.
\]

L'exposant est négatif car le nombre initial est très inférieur à $1$.

\section{11. Méthode}

\begin{methode}
\begin{enumerate}
\item déplacer la virgule pour obtenir un coefficient entre $1$ et $10$ en valeur absolue ;
\item compter les déplacements ;
\item choisir le signe de l'exposant ;
\item vérifier en revenant à l'écriture décimale.
\end{enumerate}
\end{methode}

\section{12. Revenir à l'écriture décimale}

\[
6{,}02\times10^4=60\,200.
\]

Et :
\[
6{,}02\times10^{-4}=0{,}000602.
\]

\section{13. Comparer}

Entre :
\[
7{,}3\times10^8
\]
et :
\[
4{,}9\times10^9,
\]
le second est plus grand car son exposant de $10$ est supérieur.

À exposants identiques, on compare les coefficients.

\section{14. Ordre de grandeur}

La notation scientifique rend visibles les échelles.

Le nombre :
\[
6{,}4\times10^7
\]
est de l'ordre de quelques dizaines de millions.

Selon le contexte, on peut l'associer à la puissance de dix :
\[
10^8.
\]

\section{15. Produit de puissances}

À partir de la définition :
\[
10^3\times10^4
\]
contient au total sept facteurs $10$.

Donc :
\[
10^3\times10^4=10^7.
\]

\begin{remarque}
Le programme n'exige pas la mémorisation d'un catalogue de formules générales : il faut savoir retrouver les règles par la définition.
\end{remarque}

\section{16. Quotient de puissances}

\[
\dfrac{10^7}{10^3}
=
10^4.
\]

On peut le voir en écrivant les facteurs puis en simplifiant les trois facteurs communs.

\section{17. Calcul scientifique}

\[
(3\times10^5)\times(2\times10^3)
=
6\times10^8.
\]

Le résultat est déjà sous forme scientifique.

\section{18. Normaliser le coefficient}

\[
(8\times10^6)\times(5\times10^2)
=
40\times10^8.
\]

On transforme :
\[
40=4\times10.
\]

Donc :
\[
40\times10^8=4\times10^9.
\]

\section{19. Base quelconque}

\[
2^5=32.
\]

Pour une base négative :
\[
(-3)^4=81.
\]

Mais :
\[
-3^4=-81.
\]

Les parenthèses changent la base de la puissance.

\section{20. Parité de l'exposant}

\[
(-2)^3=-8,
\]
\[
(-2)^4=16.
\]

Un nombre pair de facteurs négatifs donne un produit positif.

\section{21. Exemple scientifique}

Une distance :
\[
150\,000\,000\ \text{km}
\]
s'écrit :
\[
1{,}5\times10^8\ \text{km}.
\]

Une longueur :
\[
0{,}00001\ \text{m}
\]
s'écrit :
\[
1\times10^{-5}\ \text{m}.
\]

\section{22. Contrôle}

Pour :
\[
0{,}00042,
\]
un exposant positif serait incohérent.

On doit obtenir :
\[
4{,}2\times10^{-4}.
\]

\section{23. Erreurs fréquentes}

\begin{itemize}
\item Confondre $10^{-3}$ et $-10^3$.
\item Écrire un coefficient scientifique supérieur ou égal à $10$.
\item Oublier les parenthèses autour d'une base négative.
\item Additionner les exposants lors d'une addition ordinaire.
\item Déplacer la virgule dans le mauvais sens.
\item Donner un ordre de grandeur incompatible avec le nombre.
\end{itemize}

\section{À retenir}

\begin{aretenir}
Pour $n>0$ :
\[
\boxed{10^{-n}=\dfrac1{10^n}}.
\]

La notation scientifique s'écrit :
\[
\boxed{a\times10^n,\qquad1\le|a|<10}.
\]

Les puissances de dix permettent de gérer les préfixes d'unités, les très grands nombres, les très petits nombres et les ordres de grandeur.

Les règles de calcul doivent pouvoir être retrouvées à partir de la définition des puissances.
\end{aretenir}
